\documentclass[11pt,a4paper]{article}
\usepackage[spanish,es-noshorthands]{babel}
\usepackage{iftex}
\ifPDFTeX
\usepackage[utf8]{inputenc}\usepackage[T1]{fontenc}\usepackage{lmodern}
\else
\usepackage{fontspec}
\setmainfont{lmroman12-regular.otf}[BoldFont=lmroman12-bold.otf,ItalicFont=lmroman12-italic.otf,BoldItalicFont=lmroman10-bolditalic.otf]
\fi
\usepackage[margin=2.5cm]{geometry}
\usepackage{array,tabularx}
\usepackage{hyperref}
\hypersetup{hidelinks,pdftitle={Reporte de participación PC2}}
\setlength{\parindent}{0pt}
\setlength{\parskip}{5pt}
\setlength{\tabcolsep}{4pt}
\renewcommand{\arraystretch}{1.08}
\newcolumntype{L}[1]{>{\raggedright\arraybackslash}p{#1}}
\newcolumntype{Y}{>{\raggedright\arraybackslash}X}
\newcommand{\nombreGrupo}{Equipo K-means NYC TLC}
\begin{document}
\begin{center}
\textbf{1ACC0065 PROGRAMACIÓN CONCURRENTE Y DISTRIBUIDA}\\[3pt]
\textbf{REPORTE DE PARTICIPACIÓN}\\[3pt]
\textbf{CIENCIAS DE LA COMPUTACIÓN}
\end{center}

\begin{tabularx}{\textwidth}{|L{5cm}|Y|}\hline
Entrega & PC2 \\\hline
Nombre del grupo & \nombreGrupo \\\hline
Nombre de la solución & K-means concurrente sobre viajes de taxi de Nueva York \\\hline
Nombre completo Team Leader & Rody Sebastian Vilchez Marin \\\hline
\end{tabularx}

\vspace{6pt}
\begin{tabularx}{\textwidth}{|Y|}\hline
\textbf{Comentario del team leader} \\\hline
Los tres integrantes cumplieron sus tareas en el plazo. Dayana documentó la correspondencia entre el modelo Promela y la implementación en Go; Julio amplió los tests de casos borde de la versión concurrente y preparó este reporte y el enlace a los datos. La implementación, la verificación con Spin, las mediciones y el informe estuvieron a cargo del coordinador, que revisó e integró los aportes por pull request. \\\hline
\end{tabularx}

\vspace{6pt}
{\fontsize{9}{10}\selectfont
\begin{tabularx}{\textwidth}{|L{0.55cm}|L{2.65cm}|Y|L{1.25cm}|L{2.25cm}|L{1.2cm}|}\hline
\multicolumn{6}{|l|}{\textbf{Tabla de Participación} Marcar con un X de acuerdo al nivel de cumplimiento} \\\hline
{\scriptsize\textbf{Nro}} & {\scriptsize\textbf{Estudiante}} & {\scriptsize\textbf{Responsabilidades}} & {\scriptsize\textbf{Cumplió}} & {\scriptsize\textbf{Cumplimiento parcial}} & {\scriptsize\textbf{No cumplió}} \\\hline
1 & Rody Sebastian Vilchez Marin & Diseño, lectura de las seis features, K-means secuencial e inicialización k-means++; versión concurrente con workers persistentes y reducción reproducible; CLI y validación de parámetros (PR \#21). & X & & \\\hline
 & & Modelo Promela y mutante; verificación de cobertura y sincronización con Spin, generación de evidencias y explicación de los límites del modelo (PR \#21). & X & & \\\hline
 & & Benchmark reproducible; speedup, escalabilidad fuerte y débil, tamaño de bloque, memoria, CPU y punto de equilibrio. Mediciones en laptop, máquina virtual y Pixel 9a; contraste GPU/CPU y análisis de resultados (PR \#21, \#27, \#28, \#30 y \#31). & X & & \\\hline
 & & Informe PC2, tablas desde JSON, capturas, historial y correcciones de plataformas y participación. Actualización de README y autores; mantenimiento de CI y notebooks; atención a revisiones e integración de PR del equipo (PR \#9, \#19--\#26, \#29, \#32 y \#36). & X & & \\\hline
2 & Dayana Kety Gómez Rodriguez & Documentar la correspondencia entre Go y Promela: canal, workers, acumuladores y barrera (056a05f). & X & & \\\hline
 & & Explicar una ronda, las simplificaciones y las propiedades verificadas (056a05f). & X & & \\\hline
 & & Aclarar resultados de Spin y contraejemplo del mutante; actualizar su participación (056a05f, 44e2e83). & X & & \\\hline
 & & Corregir la documentación e incorporar los cambios mediante PR \#35 (d9c6ba2). & X & & \\\hline
3 & Julio Cesar Meza Alfaro & Probar que un cluster vacío conserva su centroide (5fbcf6b). & X & & \\\hline
 & & Probar que un empate de distancia elige el centroide de menor índice (5fbcf6b). & X & & \\\hline
 & & Probar la media con $k=1$ y la inercia cero con $k=N$ (32e169b). & X & & \\\hline
 & & Ejercitar los casos con 1, 2 y 4 workers y bloques de un punto; entregar las pruebas mediante PR \#34 (5fbcf6b, 32e169b). & X & & \\\hline
\end{tabularx}}

\newpage
\textbf{CRITERIOS DE DESCUENTO POR PARTICIPACIÓN INDIVIDUAL}

La calificación del trabajo grupal se ajustará individualmente según el nivel de compromiso y cumplimiento de las tareas asignadas. Es responsabilidad del equipo reportar objetivamente el desempeño de sus integrantes en el formato de responsabilidad.

{\small
\begin{tabularx}{\textwidth}{|L{2.3cm}|Y|L{2.4cm}|L{4.2cm}|}\hline
\textbf{Nivel de cumplimiento} & \textbf{Descripción} & \textbf{Descuento} & \textbf{Justificación} \\\hline
Cumplió & El estudiante entregó todas sus tareas asignadas (código, documentación, video) en el plazo acordado y con la calidad esperada. Participó activamente en las reuniones y la integración final. & 0 puntos (Nota completa) & Participación responsable, proactiva y organizada que contribuyó al éxito del proyecto. \\\hline
Cumplimiento Parcial & El estudiante entregó sus tareas tarde o incompletas, obligando al resto del equipo a corregir o completar su parte. Su asistencia a reuniones fue irregular. & $-5$ puntos & La falta de prolijidad o puntualidad generó una carga de trabajo adicional para el resto del equipo. \\\hline
No Cumplió & El estudiante no entregó ninguna tarea asignada, o su aporte fue irrelevante/insuficiente. No participó en la grabación del video ni en la redacción del informe. & Nota Individual: 00 (Se anula la nota del trabajo) & La inacción total perjudicó gravemente el avance del proyecto. No existe evidencia de aprendizaje ni trabajo realizado. \\\hline
\end{tabularx}}
\end{document}
